\begin{figure}[htbp]
\centering
\begin{tikzpicture}

% ========== Bisection Method ==========
\begin{axis}[
    name=bisection,
    width=7.5cm, height=7cm,
    xlabel={$x$},
    ylabel={$f(x)$},
    xmin=-0.5, xmax=4,
    ymin=-3, ymax=5,
    grid=major,
    grid style={UofSC50Black!30},
    axis lines=middle,
    xtick={0,1,2,3},
    ytick={-2,0,2,4},
    title={\textbf{Bisection Method}},
    title style={at={(0.5,1.02)}, font=\bfseries},
    clip=false
]

% Function: f(x) = x^2 - 3 (root at sqrt(3) ~ 1.732)
\addplot[UofSCBlack, very thick, domain=-0.3:3.5, samples=100] {x^2 - 3};

% Root location
\pgfmathsetmacro{\rootval}{sqrt(3)}

% Initial interval [a, b] = [1, 3]
\pgfmathsetmacro{\aone}{1}
\pgfmathsetmacro{\bone}{3}
\pgfmathsetmacro{\cone}{(\aone + \bone)/2} % = 2

% Draw interval brackets - Iteration 1
\draw[UofSCGarnet, very thick] (axis cs:\aone,-2.8) -- (axis cs:\aone,-2.5);
\draw[UofSCGarnet, very thick] (axis cs:\bone,-2.8) -- (axis cs:\bone,-2.5);
\draw[UofSCGarnet, thick] (axis cs:\aone,-2.65) -- (axis cs:\bone,-2.65);
\node[UofSCGarnet, font=\tiny, below] at (axis cs:{(\aone+\bone)/2},-2.8) {$[a_0, b_0]$};

% Midpoint c1 = 2
\draw[UofSCGarnet, dashed, thick] (axis cs:\cone,-2.5) -- (axis cs:\cone,{(\cone)^2-3});
\node[circle, fill=UofSCGarnet, minimum size=4pt, inner sep=0pt] at (axis cs:\cone,{(\cone)^2-3}) {};

% Iteration 2: f(2) = 1 > 0, so new interval [1, 2]
\pgfmathsetmacro{\atwo}{1}
\pgfmathsetmacro{\btwo}{2}
\pgfmathsetmacro{\ctwo}{(\atwo + \btwo)/2} % = 1.5

\draw[UofSCAtlantic, very thick] (axis cs:\atwo,-2.1) -- (axis cs:\atwo,-1.8);
\draw[UofSCAtlantic, very thick] (axis cs:\btwo,-2.1) -- (axis cs:\btwo,-1.8);
\draw[UofSCAtlantic, thick] (axis cs:\atwo,-1.95) -- (axis cs:\btwo,-1.95);
\node[UofSCAtlantic, font=\tiny, below] at (axis cs:{(\atwo+\btwo)/2},-2.1) {$[a_1, b_1]$};

\draw[UofSCAtlantic, dashed, thick] (axis cs:\ctwo,-1.8) -- (axis cs:\ctwo,{(\ctwo)^2-3});
\node[circle, fill=UofSCAtlantic, minimum size=4pt, inner sep=0pt] at (axis cs:\ctwo,{(\ctwo)^2-3}) {};

% Iteration 3: f(1.5) = -0.75 < 0, so new interval [1.5, 2]
\pgfmathsetmacro{\athree}{1.5}
\pgfmathsetmacro{\bthree}{2}
\pgfmathsetmacro{\cthree}{(\athree + \bthree)/2} % = 1.75

\draw[UofSCHorseshoe, very thick] (axis cs:\athree,-1.4) -- (axis cs:\athree,-1.1);
\draw[UofSCHorseshoe, very thick] (axis cs:\bthree,-1.4) -- (axis cs:\bthree,-1.1);
\draw[UofSCHorseshoe, thick] (axis cs:\athree,-1.25) -- (axis cs:\bthree,-1.25);
\node[UofSCHorseshoe, font=\tiny, below] at (axis cs:{(\athree+\bthree)/2},-1.4) {$[a_2, b_2]$};

\draw[UofSCHorseshoe, dashed, thick] (axis cs:\cthree,-1.1) -- (axis cs:\cthree,{(\cthree)^2-3});
\node[circle, fill=UofSCHorseshoe, minimum size=4pt, inner sep=0pt] at (axis cs:\cthree,{(\cthree)^2-3}) {};

% Mark the actual root
\node[circle, draw=UofSCBlack, fill=white, minimum size=6pt, inner sep=0pt, thick] at (axis cs:\rootval,0) {};
\node[UofSCBlack, font=\small, anchor=south west] at (axis cs:{\rootval+0.1},0.2) {$x^* = \sqrt{3}$};

% Sign indicators
\node[UofSCGarnet, font=\footnotesize] at (axis cs:0.7,-1.8) {$f(a)<0$};
\node[UofSCGarnet, font=\footnotesize] at (axis cs:3.2,5.5) {$f(b)>0$};

\end{axis}

% ========== Newton-Raphson Method ==========
\begin{axis}[
    name=newton,
    at={(bisection.south east)},
    anchor=south west,
    xshift=1.5cm,
    width=7.5cm, height=7cm,
    xlabel={$x$},
    ylabel={$f(x)$},
    xmin=-0.5, xmax=4,
    ymin=-3, ymax=5,
    grid=major,
    grid style={UofSC50Black!30},
    axis lines=middle,
    xtick={0,1,2,3},
    ytick={-2,0,2,4},
    title={\textbf{Newton-Raphson Method}},
    title style={at={(0.5,1.02)}, font=\bfseries},
    clip=false
]

% Function: f(x) = x^2 - 3
\addplot[UofSCBlack, very thick, domain=-0.3:3.5, samples=100] {x^2 - 3};

% Root location
\pgfmathsetmacro{\rootval}{sqrt(3)}

% Newton iteration: x_{n+1} = x_n - f(x_n)/f'(x_n) = x_n - (x_n^2 - 3)/(2*x_n)
% Starting from x_0 = 3
\pgfmathsetmacro{\xzero}{3}
\pgfmathsetmacro{\fzero}{(\xzero)^2 - 3} % = 6
\pgfmathsetmacro{\fpzero}{2*\xzero} % = 6
\pgfmathsetmacro{\xone}{\xzero - \fzero/\fpzero} % = 3 - 6/6 = 2

% Point and tangent at x_0 = 3
\node[circle, fill=UofSCGarnet, minimum size=5pt, inner sep=0pt] at (axis cs:\xzero,\fzero) {};
\addplot[UofSCGarnet, thick, domain=1.5:3.5, samples=2, forget plot] {\fpzero*(x - \xzero) + \fzero};
\node[UofSCGarnet, font=\small, anchor=south] at (axis cs:\xzero,{\fzero+0.3}) {$(x_0, f(x_0))$};

% x_1 = 2
\draw[UofSCGarnet, dashed] (axis cs:\xone,0) -- (axis cs:\xone,{(\xone)^2 - 3});
\node[UofSCGarnet, font=\footnotesize, below] at (axis cs:\xone,-0.3) {$x_1$};

% Iteration 2: from x_1 = 2
\pgfmathsetmacro{\fone}{(\xone)^2 - 3} % = 1
\pgfmathsetmacro{\fpone}{2*\xone} % = 4
\pgfmathsetmacro{\xtwo}{\xone - \fone/\fpone} % = 2 - 1/4 = 1.75

\node[circle, fill=UofSCAtlantic, minimum size=5pt, inner sep=0pt] at (axis cs:\xone,\fone) {};
\addplot[UofSCAtlantic, thick, domain=1.2:2.5, samples=2, forget plot] {\fpone*(x - \xone) + \fone};
\node[UofSCAtlantic, font=\small, anchor=south west] at (axis cs:\xone,{\fone+0.2}) {$(x_1, f(x_1))$};

% x_2 = 1.75
\draw[UofSCAtlantic, dashed] (axis cs:\xtwo,0) -- (axis cs:\xtwo,{(\xtwo)^2 - 3});
\node[UofSCAtlantic, font=\footnotesize, below] at (axis cs:\xtwo,-0.3) {$x_2$};

% Iteration 3: from x_2 = 1.75
\pgfmathsetmacro{\ftwo}{(\xtwo)^2 - 3} % = 0.0625
\pgfmathsetmacro{\fptwo}{2*\xtwo} % = 3.5
\pgfmathsetmacro{\xthree}{\xtwo - \ftwo/\fptwo} % ~ 1.732143

\node[circle, fill=UofSCHorseshoe, minimum size=5pt, inner sep=0pt] at (axis cs:\xtwo,\ftwo) {};
\addplot[UofSCHorseshoe, thick, domain=1.4:2.1, samples=2, forget plot] {\fptwo*(x - \xtwo) + \ftwo};

% x_3 ~ 1.732 (very close to root)
\node[UofSCHorseshoe, font=\footnotesize, below] at (axis cs:\xthree,-0.3) {$x_3$};

% Mark the actual root
\node[circle, draw=UofSCBlack, fill=white, minimum size=6pt, inner sep=0pt, thick] at (axis cs:\rootval,0) {};

% Tangent line concept
\node[align=left, font=\footnotesize, anchor=north west] at (axis cs:2.2,4.8) {
    Tangent at $(x_n, f(x_n))$:\\
    slope $= f'(x_n)$
};

\end{axis}

% ========== Convergence comparison ==========
\begin{axis}[
    name=convergence,
    at={(bisection.south west)},
    anchor=north west,
    yshift=-1.5cm,
    width=14.5cm, height=5cm,
    xlabel={Iteration $n$},
    ylabel={$|x_n - x^*|$ (log scale)},
    xmin=0, xmax=10,
    ymin=1e-10, ymax=10,
    ymode=log,
    grid=major,
    grid style={UofSC50Black!30},
    title={\textbf{Convergence Comparison}},
    title style={at={(0.5,1.02)}, font=\bfseries},
    legend pos=north east,
    legend style={font=\footnotesize, sharp corners, draw=UofSC70Black, fill=white}
]

% Bisection: linear convergence, error halves each iteration
% Starting interval [1, 3], root = sqrt(3) ~ 1.732
% Error roughly halves: 1, 0.5, 0.25, 0.125, ...
\addplot[UofSCGarnet, very thick, mark=*, mark size=2.5pt, mark options={fill=UofSCGarnet}]
    coordinates {
        (0, 1.268) (1, 0.268) (2, 0.232) (3, 0.018) (4, 0.107)
        (5, 0.045) (6, 0.014) (7, 0.0155) (8, 0.00078) (9, 0.00722)
        (10, 0.00322)
    };
\addlegendentry{Bisection (linear)}

% Newton: quadratic convergence
% x_0 = 3, error = 1.268
% x_1 = 2, error = 0.268
% x_2 = 1.75, error = 0.018
% x_3 = 1.732143, error = 0.000143
% x_4 ~ root, error ~ 1e-8
\addplot[UofSCAtlantic, very thick, mark=square*, mark size=2.5pt, mark options={fill=UofSCAtlantic}]
    coordinates {
        (0, 1.268) (1, 0.268) (2, 0.018) (3, 0.000111)
        (4, 4.2e-9) (5, 1e-10)
    };
\addlegendentry{Newton (quadratic)}

% Secant: superlinear convergence (order ~ 1.618)
\addplot[UofSCHorseshoe, very thick, mark=triangle*, mark size=2.5pt, mark options={fill=UofSCHorseshoe}]
    coordinates {
        (0, 1.268) (1, 0.732) (2, 0.134) (3, 0.0193)
        (4, 0.00066) (5, 3.3e-6) (6, 5.6e-10)
    };
\addlegendentry{Secant (superlinear)}

% Reference lines for convergence rates
\draw[UofSC50Black, thick, dashed] (axis cs:0,1) -- (axis cs:10,0.001);
\node[UofSC50Black, font=\tiny, anchor=west, rotate=-13] at (axis cs:6,0.015) {Linear $O(1/2^n)$};

\end{axis}

% ========== Method formulas ==========
\node[anchor=north west, align=left, font=\small] at (0,-10.5) {
    \textcolor{UofSCGarnet}{\textbf{Bisection:}} $c = \frac{a+b}{2}$; if $f(a)f(c) < 0$: $b \leftarrow c$, else $a \leftarrow c$ \quad (Linear: error halves each step)
};
\node[anchor=north west, align=left, font=\small] at (0,-11.2) {
    \textcolor{UofSCAtlantic}{\textbf{Newton-Raphson:}} $x_{n+1} = x_n - \dfrac{f(x_n)}{f'(x_n)}$ \quad (Quadratic: digits of accuracy double each step)
};
\node[anchor=north west, align=left, font=\small] at (0,-12.0) {
    \textcolor{UofSCHorseshoe}{\textbf{Secant:}} $x_{n+1} = x_n - f(x_n)\dfrac{x_n - x_{n-1}}{f(x_n) - f(x_{n-1})}$ \quad (Superlinear: order $\approx 1.618$)
};

\end{tikzpicture}
\caption{Root finding methods for $f(x) = x^2 - 3$ with root $x^* = \sqrt{3} \approx 1.732$. \textbf{Top left}: Bisection method iteratively halves the bracketing interval $[a,b]$ where $f(a)f(b) < 0$. It is guaranteed to converge but only linearly. \textbf{Top right}: Newton-Raphson method follows tangent lines from the current approximation to the $x$-axis, achieving quadratic convergence when sufficiently close to the root. \textbf{Bottom}: Convergence comparison showing that Newton's method (Atlantic blue) converges in about 5 iterations to machine precision, while bisection (Garnet) requires many more iterations for the same accuracy. The secant method (Horseshoe green) achieves superlinear convergence without requiring derivative evaluation.}
\label{fig:root_finding}
\end{figure}
